\documentclass{article}
\usepackage[T1]{fontenc}
\usepackage{lmodern}
\usepackage{graphicx}
\usepackage{amsmath,amssymb}
\usepackage[a4paper,margin=2cm]{geometry}
\usepackage[absolute,overlay]{textpos}
\setlength{\TPHorizModule}{1mm}
\setlength{\TPVertModule}{1mm}
\usepackage[indonesian]{babel}
\usepackage{hyperref}
\setlength{\emergencystretch}{2em}
% unit: Halaman 15

\begin{document}

% === Halaman 15 (llm) ===

\section*{2. Serangan berbasis statistik (statistical attack)}

\begin{itemize}
    \item Menggunakan analisis matematik pada citra untuk menemukan perbedaan antara \textit{cover image} dengan \textit{stego image}.
    \item Didasarkan pada fakta bahwa penyembunyian pesan ke dalam media menimbulkan artefak yang dapat dideteksi secara statistik sehingga dapat mengungkap penyembunyian pesan atau pesan yang disembunyikan itu sendiri.
    \item Contoh serangan statistik:
    \begin{enumerate}[a.]
        \item \textit{histogram analysis}
        \item \textit{Regular-singular (RS) analysis}
        \item \textit{Chi-square analysis}
        \item \textit{Sample pair (SP) analysis}
    \end{enumerate}
\end{itemize}

\end{document}
